
\documentclass{article}
\usepackage{colortbl}
\usepackage{makecell}
\usepackage{multirow}
\usepackage{supertabular}

\begin{document}

\newcounter{utterance}

\centering \large Interaction Transcript for game `imagegame', experiment `compact\_grids', episode 10 with reka{-}flash{-}3.1{-}t0.0.
\vspace{24pt}

{ \footnotesize  \setcounter{utterance}{1}
\setlength{\tabcolsep}{0pt}
\begin{supertabular}{c@{$\;$}|p{.15\linewidth}@{}p{.15\linewidth}p{.15\linewidth}p{.15\linewidth}p{.15\linewidth}p{.15\linewidth}}
    \# & $\;$A & \multicolumn{4}{c}{Game Master} & $\;\:$B\\
    \hline

    \theutterance \stepcounter{utterance}  
    & & \multicolumn{4}{p{0.6\linewidth}}{
        \cellcolor[rgb]{0.9,0.9,0.9}{
            \makecell[{{p{\linewidth}}}]{
                \texttt{\tiny{[P1$\langle$GM]}}
                \texttt{Let us do the following task. The goal is to fill an empty grid that looks like this:} \\
\\ 
\texttt{▢ ▢ ▢ ▢ ▢} \\
\texttt{▢ ▢ ▢ ▢ ▢} \\
\texttt{▢ ▢ ▢ ▢ ▢} \\
\texttt{▢ ▢ ▢ ▢ ▢} \\
\texttt{▢ ▢ ▢ ▢ ▢} \\
\\ 
\texttt{A filled grid below is 5 by 5 and can look like this:} \\
\\ 
\texttt{▢ ▢ ▢ ▢ ▢} \\
\texttt{▢ ▢ E ▢ ▢} \\
\texttt{▢ ▢ ▢ ▢ ▢} \\
\texttt{▢ ▢ ▢ ▢ ▢} \\
\texttt{X X X X X} \\
\\ 
\texttt{I want you to describe this grid to me, step by step. You don't need to describe the empty squares, which are denoted with "▢". Only describe the location of letters in the grid. Then you wait for me to say "What is your next command?", and then you continue with the next step. Take the size of the grid into consideration while giving commands. You can give only one command at each turn but it can refer to single or multiple cells, rows, columns or even the whole grid at once. When you have described everything, you say "Command: DONE". Only give the command and omit everything else.} \\
\\ 
\texttt{For the filled grid above, here are the example steps.} \\
\\ 
\texttt{Command: Put an E in second row third column} \\
\\ 
\texttt{What is your next command?} \\
\texttt{Command: Fill the last row with X} \\
\\ 
\texttt{What is your next command?} \\
\texttt{Command: DONE} \\
\\ 
\texttt{Ok. Please do this for the following example, which is a 5 by 5 grid.} \\
\texttt{Q Q ▢ Q Q} \\
\texttt{Q ▢ ▢ ▢ Q} \\
\texttt{▢ ▢ ▢ ▢ ▢} \\
\texttt{Q ▢ ▢ ▢ Q} \\
\texttt{Q Q ▢ Q Q} \\
\texttt{What is your next command?} \\
            }
        }
    }
    & \\ \\

    \theutterance \stepcounter{utterance}  
    & \multicolumn{4}{p{0.6\linewidth}}{
        \cellcolor[rgb]{0.9,0.9,0.9}{
            \makecell[{{p{\linewidth}}}]{
                \texttt{\tiny{[P1$\rangle$GM]}}
                \texttt{<reasoning>} \\
\texttt{Let's see. The user provided a 5x5 grid filled with some Qs and a blank, and they want me to describe the steps to fill it, only mentioning the non{-}blank cells. The example they gave earlier had commands like putting an E in a specific cell and filling the last row with Xs.} \\
\\ 
\texttt{First, I need to look at the given grid:} \\
\\ 
\texttt{Row 1: Q Q ▢ Q Q} \\
\texttt{Row 2: Q ▢ ▢ ▢ Q} \\
\texttt{Row 3: ▢ ▢ ▢ ▢ ▢} \\
\texttt{Row 4: Q ▢ ▢ ▢ Q} \\
\texttt{Row 5: Q Q ▢ Q Q} \\
\\ 
\texttt{So, the empty cells are the third column in each row. Wait, let me check each row:} \\
\\ 
\texttt{Row 1: Columns 1,2,4,5 are Q; column 3 is ▢.} \\
\texttt{Row 2: Columns 1 and 5 are Q; columns 2{-}4 are ▢.} \\
\texttt{Row 3: All are ▢.} \\
\texttt{Row 4: Columns 1 and 5 are Q; columns 2{-}4 are ▢.} \\
\texttt{Row 5: Columns 1,2,4,5 are Q; column 3 is ▢.} \\
\\ 
\texttt{So the filled cells are the Qs. The task is to describe the locations of the Qs. The user wants commands that refer to these positions. Each command can be placing a Q in specific cells, filling entire rows or columns, etc.} \\
\\ 
\texttt{Looking at the example, they used "Put an E in second row third column" and "Fill the last row with X". So similar commands here. Since all the non{-}blank cells are Qs, the commands should specify where the Qs are placed.} \\
\\ 
\texttt{First, notice that in rows 1, 2, 4, 5, the third column is empty. In row 3, all are empty. So the Qs are in positions where the grid isn't blank. So the steps would involve placing Qs in the non{-}blank cells.} \\
\\ 
\texttt{But how to group them? Let's see:} \\
\\ 
\texttt{Looking at the grid:} \\
\\ 
\texttt{Row 1: Qs in columns 1,2,4,5. So four Qs here. The third is empty.} \\
\\ 
\texttt{Row 2: Qs in 1 and 5. Columns 2{-}4 are empty.} \\
\\ 
\texttt{Row 3: All empty.} \\
            }
        }
    }
    & & \\ \\

    \theutterance \stepcounter{utterance}  
    & & & \multicolumn{2}{p{0.3\linewidth}}{
        \cellcolor[rgb]{0.9,0.9,0.9}{
            \makecell[{{p{\linewidth}}}]{
                \texttt{\tiny{[GM$|$GM]}}
                \texttt{Invalid instruction format} \\
            }
        }
    }
    & & \\ \\

\end{supertabular}
}

\end{document}
